% Source: https://artofproblemsolving.com/community/c3h1936254_how_does_the_triangle_inequality_work
% Stable identity: aops-topic-1936254-q01
(a) Let triangle $ABC$ have sides $a=BC$, $b=CA$, $c=AB$, and suppose $b$ is a longest side. Then $\angle B$ is a largest angle. If $\angle A$ or $\angle C$ were at least $90^\circ$, then $\angle B$ would be too, and the angle sum would exceed $180^\circ$. So the angles at $A$ and $C$ are acute, and the foot $H$ of the altitude from $B$ lies strictly between $A$ and $C$.

\begin{asy}
usepackage("amsmath");
size(6cm);
defaultpen(fontsize(10pt));
pair A=(0,0),C=(7,0),B=(2.2,2.6),H=(2.2,0);
draw(A--B--C--cycle,linewidth(1));
draw(B--H,dashed);draw(H+(0.25,0)--H+(0.25,0.25)--H+(0,0.25));
label("$A$",A,SW);label("$C$",C,SE);label("$B$",B,N);label("$H$",H,S);
label("$c$",(A+B)/2,NW);label("$a$",(B+C)/2,NE);label("$b$",(A+C)/2+(0.8,0),S);
label("\mbox{}",(min(currentpicture,user=true)+max(currentpicture,user=true))/2,black); // keeps the text after the figure black when LaTeX typesets the labels inline
\end{asy}

In the right triangles $ABH$ and $CBH$ the hypotenuses are longer than the legs, because $BH>0$: $c>AH$ and $a>HC$. Adding, $a+c>AH+HC=b$. The other two inequalities $a+b>c$ and $b+c>a$ are immediate, since $b$ is at least as long as $a$ and $c$ and all sides are positive. So each side is less than the sum of the other two, with strict inequality.

(b) For three positive lengths it is enough to check that the longest is less than the sum of the other two. $(3,4,5)$: $3+4>5$, so it is a triangle. $(2,3,6)$: $2+3<6$, not a triangle. $(2,3,5)$: $2+3=5$, which is not a triangle: by (a) a genuine triangle needs strict inequality, and lengths with $AB+BC=AC$ force $B$ to lie on segment $AC$, a ``flattened'' triangle. So only $\boxed{(3,4,5)}$ works.
